\clearpage
\appendix
\section{AI-credit lexicon}\label{app:lexicon}
Patterns are matched case-insensitively against the concatenated credit strings of each CVE record after removing the exclusion patterns (a human researcher whose given name is Claude; a username containing ``gpt''; the brand name of a human bug-bounty programme containing the letters ``AI''). When several patterns match, the lowest tier wins.
{\small
\begin{longtable}{p{0.24\textwidth}p{0.13\textwidth}p{0.04\textwidth}p{0.45\textwidth}}
\caption{AI-credit lexicon (case-insensitive regular expressions applied to CVE credit text).}\label{tab:lexicon}\\
\toprule
Label & Developer & Tier & Pattern (truncated) \\
\midrule
\endfirsthead
\toprule
Label & Developer & Tier & Pattern (truncated) \\
\midrule
\endhead
\bottomrule
\endfoot
Big Sleep & Google & 1 & \texttt{big ?sleep} \\
XBOW & XBOW & 1 & \texttt{\textbackslash{}bxbow\textbackslash{}b|\allowbreak{}xbow\textbackslash{}.com} \\
Claude / Project Glasswing & Anthropic & 1 & \texttt{\textbackslash{}bclaude\textbackslash{}b|\allowbreak{}anthropic|\allowbreak{}glasswing|\allowbreak{}powered by mythos} \\
Codex Security / Aardvark & OpenAI & 1 & \texttt{\textbackslash{}bcodex\textbackslash{}b|\allowbreak{}aardvark} \\
GPT model & OpenAI & 1 & \texttt{\textbackslash{}bgpt-?\textbackslash{}d|\allowbreak{}gpt-5} \\
Other named AI tool & Various & 1 & \texttt{\textbackslash{}bAI (?:SAST|\allowbreak{}assistant|\allowbreak{}agent|\allowbreak{}agents)\textbackslash{}b|\allowbreak{}security AI agent|\allowbreak{}\textbackslash{}bAI-agent\textbackslash{}b|\allowbreak{}} \\
OpenAI (organisation) & OpenAI & 2 & \texttt{\textbackslash{}bopenai\textbackslash{}b} \\
AISLE & AISLE & 2 & \texttt{\textbackslash{}baisle\textbackslash{}b} \\
depthfirst & depthfirst & 2 & \texttt{depthfirst} \\
ZeroPath & ZeroPath & 2 & \texttt{zeropath} \\
Striga & Striga & 2 & \texttt{striga\textbackslash{}.ai} \\
ZAST.AI & ZAST.AI & 2 & \texttt{zast\textbackslash{}.ai} \\
BugBunny.ai & BugBunny.ai & 2 & \texttt{bugbunny\textbackslash{}.ai} \\
CodeAnt AI & CodeAnt AI & 2 & \texttt{codeant ai} \\
Hacktron & Hacktron & 2 & \texttt{hacktron} \\
Xint & Theori & 2 & \texttt{\textbackslash{}bxint\textbackslash{}b} \\
WatchGuard AI Security Research & WatchGuard & 2 & \texttt{watchguard ai security research} \\
Mopmonk AI & Mopmonk AI & 2 & \texttt{mopmonk} \\
Z.ai & Zhipu AI & 2 & \texttt{\textbackslash{}bz\textbackslash{}.ai\textbackslash{}b} \\
Moonshot AI & Moonshot AI & 2 & \texttt{moonshot\textbackslash{}.ai} \\
Horizon3.ai & Horizon3.ai & 2 & \texttt{horizon3\textbackslash{}.ai} \\
Other AI-native organisation & Various & 2 & \texttt{pwn\textbackslash{}.ai|\allowbreak{}vexera|\allowbreak{}submersion|\allowbreak{}aether ai|\allowbreak{}innora\textbackslash{}.ai|\allowbreak{}sectum ai|\allowbreak{}parameter\textbackslash{}.ai} \\
\end{longtable}
}

\section{KEV category rules and banking-relevance weights}\label{app:rules}
{\small
\begin{longtable}{p{0.04\textwidth}p{0.26\textwidth}p{0.58\textwidth}}
\caption{KEV product-category rules (ordered; first match on vendor and product, then on vulnerability name). Patterns truncated; full rules in \texttt{data/curated/kev\_category\_rules.csv}.}\label{tab:rules}\\
\toprule
\# & Category & Pattern (truncated) \\
\midrule
\endfirsthead
\toprule
\# & Category & Pattern (truncated) \\
\midrule
\endhead
\bottomrule
\endfoot
1 & AI \& automation platforms & \texttt{langflow|\allowbreak{}litellm|\allowbreak{}ollama|\allowbreak{}\textbackslash{}bray\textbackslash{}b|\allowbreak{}mlflow|\allowbreak{}comfyui|\allowbreak{}anything-?llm|\allowbreak{}open-?webui|\allowbreak{}langcha\ldots} \\
2 & Managed file transfer & \texttt{moveit|\allowbreak{}goanywhere|\allowbreak{}accellion|\allowbreak{}\textbackslash{}bfta\textbackslash{}b|\allowbreak{}cleo|\allowbreak{}crushftp|\allowbreak{}serv-u|\allowbreak{}wing ?ftp|\allowbreak{}kiteworks|\allowbreak{}glo\ldots} \\
3 & Email \& collaboration & \texttt{exchange server|\allowbreak{}exchange|\allowbreak{}sharepoint|\allowbreak{}zimbra|\allowbreak{}roundcube|\allowbreak{}smartermail|\allowbreak{}\textbackslash{}bexim\textbackslash{}b|\allowbreak{}conflu\ldots} \\
4 & IT management \& remote access tools & \texttt{screenconnect|\allowbreak{}kaseya|\allowbreak{}simplehelp|\allowbreak{}beyondtrust|\allowbreak{}manageengine|\allowbreak{}solarwinds|\allowbreak{}n-able|\allowbreak{}nagio\ldots} \\
5 & Virtualisation \& backup & \texttt{vmware|\allowbreak{}vcenter|\allowbreak{}esxi|\allowbreak{}vsphere|\allowbreak{}veeam|\allowbreak{}veritas|\allowbreak{}backup|\allowbreak{}hyper-v|\allowbreak{}xenserver|\allowbreak{}proxmox|\allowbreak{}nakiv\ldots} \\
6 & Enterprise edge \& network infrastructure & \texttt{fortinet|\allowbreak{}fortios|\allowbreak{}fortiproxy|\allowbreak{}fortiweb|\allowbreak{}fortigate|\allowbreak{}fortimanager|\allowbreak{}fortisiem|\allowbreak{}palo alto|\allowbreak{}\ldots} \\
7 & Browsers \& plug-ins & \texttt{chromium|\allowbreak{}chrome|\allowbreak{}\textbackslash{}bv8\textbackslash{}b|\allowbreak{}internet explorer|\allowbreak{}\textbackslash{}bedge\textbackslash{}b|\allowbreak{}firefox|\allowbreak{}safari|\allowbreak{}webkit|\allowbreak{}flash pl\ldots} \\
8 & Mobile platforms & \texttt{android|\allowbreak{}\textbackslash{}bios\textbackslash{}b|\allowbreak{}ipados|\allowbreak{}watchos|\allowbreak{}samsung|\allowbreak{}pixel|\allowbreak{}multiple products.*apple|\allowbreak{}apple.*mul\ldots} \\
9 & Office \& document software & \texttt{office|\allowbreak{}\textbackslash{}bword\textbackslash{}b|\allowbreak{}excel|\allowbreak{}powerpoint|\allowbreak{}outlook|\allowbreak{}publisher|\allowbreak{}acrobat|\allowbreak{}reader|\allowbreak{}winrar|\allowbreak{}7-?zip|\allowbreak{}\ldots} \\
10 & Operating systems \& identity & \texttt{windows|\allowbreak{}win32k|\allowbreak{}\textbackslash{}bkernel\textbackslash{}b|\allowbreak{}linux|\allowbreak{}macos|\allowbreak{}os x|\allowbreak{}active directory|\allowbreak{}kerberos|\allowbreak{}ntlm|\allowbreak{}smbv?\textbackslash{}\ldots} \\
11 & Security \& endpoint software & \texttt{trend micro|\allowbreak{}defender|\allowbreak{}symantec|\allowbreak{}mcafee|\allowbreak{}kaspersky|\allowbreak{}eset|\allowbreak{}bitdefender|\allowbreak{}crowdstrike|\allowbreak{}sent\ldots} \\
12 & Hardware firmware \& chipsets & \texttt{qualcomm|\allowbreak{}chipset|\allowbreak{}\textbackslash{}barm\textbackslash{}b|\allowbreak{}mali|\allowbreak{}gpu|\allowbreak{}intel|\allowbreak{}\textbackslash{}bamd\textbackslash{}b|\allowbreak{}mediatek|\allowbreak{}gigabyte|\allowbreak{}firmware|\allowbreak{}bios|\allowbreak{}\ldots} \\
13 & SOHO routers NAS \& IoT & \texttt{d-link|\allowbreak{}dlink|\allowbreak{}netgear|\allowbreak{}tp-?link|\allowbreak{}tenda|\allowbreak{}linksys|\allowbreak{}totolink|\allowbreak{}qnap|\allowbreak{}synology|\allowbreak{}hikvision|\allowbreak{}dah\ldots} \\
14 & Industrial control \& OT & \texttt{siemens|\allowbreak{}schneider|\allowbreak{}rockwell|\allowbreak{}unitronics|\allowbreak{}moxa|\allowbreak{}\textbackslash{}bge\textbackslash{}b|\allowbreak{}general electric|\allowbreak{}scada|\allowbreak{}ics\textbackslash{}b|\allowbreak{}p\ldots} \\
15 & Enterprise applications \& web frameworks & \texttt{oracle|\allowbreak{}weblogic|\allowbreak{}e-business|\allowbreak{}peoplesoft|\allowbreak{}jd edwards|\allowbreak{}fusion middleware|\allowbreak{}\textbackslash{}bsap\textbackslash{}b|\allowbreak{}netwe\ldots} \\
\end{longtable}
\begin{longtable}{p{0.26\textwidth}p{0.06\textwidth}p{0.56\textwidth}}
\caption{Banking-relevance weights $w_c$ by product category (author judgement; sensitivity analysed in Section~\ref{sec:exploitation}).}\label{tab:weights}\\
\toprule
Category & $w_c$ & Rationale \\
\midrule
\endfirsthead
\toprule
Category & $w_c$ & Rationale \\
\midrule
\endhead
\bottomrule
\endfoot
Enterprise edge \& network infrastructure & 1.0 & Perimeter VPN/firewall/ADC appliances front staff remote access and third-party links; pre-authentication bugs give direct entry \\
Managed file transfer & 1.0 & Used for bank-to-bank batch files and corporate customer data exchange; repeated mass-exploitation campaigns \\
IT management \& remote access tools & 0.9 & Privileged tooling (RMM; help desk; CI/CD) gives fleet-wide control and is often run by outsourcers \\
Email \& collaboration & 0.9 & On-premise mail and collaboration servers hold sensitive data and staff credentials \\
Virtualisation \& backup & 0.8 & Hypervisor and backup compromise enables ransomware against core and recovery systems simultaneously \\
Enterprise applications \& web frameworks & 0.8 & ERP; middleware and web frameworks underpin internet banking back ends and internal workflow \\
Operating systems \& identity & 0.7 & Post-exploitation privilege escalation and directory compromise inside the estate \\
AI \& automation platforms & 0.7 & Newly adopted LLM gateways and workflow tools; often deployed quickly with broad data access \\
Mobile platforms & 0.6 & Customer devices running mobile banking; exploited mainly for fraud and account takeover rather than bank intrusion \\
Browsers \& plug-ins & 0.6 & Employee endpoints; initial access via drive-by or targeted links \\
Office \& document software & 0.6 & Employee endpoints; initial access via documents \\
Security \& endpoint software & 0.6 & Agents with high privilege on every endpoint \\
Hardware firmware \& chipsets & 0.4 & Relevant to HSM; ATM and server firmware but rarely the first stage \\
SOHO routers NAS \& IoT & 0.2 & Mostly consumer and small-office devices; relevant to branch and home-worker edge and botnets \\
Industrial control \& OT & 0.1 & Data-centre facilities only; marginal for banking \\
Unclassified & 0.5 & Neutral weight for unmatched products \\
\end{longtable}
}

\section{AI vulnerability-discovery milestones}\label{app:milestones}
\begin{table}[h]
\centering
\caption{Selected milestones (author-coded; verify dates and descriptions against the cited primary sources).}
\label{tab:milestones}
\scriptsize
\begin{tabular}{p{0.07\textwidth}p{0.55\textwidth}p{0.15\textwidth}p{0.13\textwidth}}
\toprule
Date & Event & Actor & Kind \\
\midrule
2016-08 & DARPA Cyber Grand Challenge: fully autonomous systems find and patch bugs in a closed environment & DARPA & Capability \\
2024-02 & Threat-intelligence reports describe state-linked actors using LLMs for reconnaissance and scripting & Microsoft and OpenAI & Misuse \\
2024-04 & Academic study reports an LLM agent exploiting most of a small set of one-day web vulnerabilities given their advisories & University researchers & Capability \\
2024-11 & Big Sleep agent reports an exploitable memory-safety bug in SQLite before release & Google & Defensive discovery \\
2024-11 & AI-generated fuzz targets in OSS-Fuzz report new vulnerabilities including one in OpenSSL & Google & Defensive discovery \\
2025-03 & First CVE in our data crediting an autonomous AI system (XBOW) & XBOW & Defensive discovery \\
2025-07 & Big Sleep credited with a SQLite flaw that the vendor says was otherwise known only to threat actors & Google & Defensive discovery \\
2025-08 & DARPA AI Cyber Challenge final: cyber-reasoning systems find and patch real open-source bugs & DARPA & Capability \\
2025-10 & Agentic security researchers announced by major labs (Aardvark; CodeMender) & OpenAI and Google DeepMind & Defensive discovery \\
2025-11 & Report of an espionage campaign in which an AI agent executed most tactical steps against about thirty organisations & Anthropic & Misuse \\
2026-04 & Restricted-access programme gives selected organisations a frontier model for defensive vulnerability discovery (Project Glasswing) & Anthropic & Defensive discovery \\
2026-10 & Monthly CVE publications exceed four times the 2025 monthly level; AI systems credited on hundreds of CVEs per month & CVE Program data & Disclosure shock \\
\bottomrule
\end{tabular}
\end{table}

\section{Curated financial-sector incidents}\label{app:incidents}
{\scriptsize
\begin{longtable}{p{0.05\textwidth}p{0.28\textwidth}p{0.40\textwidth}p{0.13\textwidth}}
\caption{Curated financial-sector incident set (purposive sample of well-documented public cases, author-coded; \textsuperscript{TH} = Thailand). Details should be verified against primary sources before citation.}\label{tab:incidents}\\
\toprule
Year & Victim & Initial vector (as publicly reported) & Zero-day? \\
\midrule
\endfirsthead
\toprule
Year & Victim & Initial vector (as publicly reported) & Zero-day? \\
\midrule
\endhead
\bottomrule
\endfoot
2011 & RSA (SecurID vendor) & Spear-phishing with Flash zero-day in spreadsheet & Yes \\
2013 & Shinhan, NongHyup and Jeju banks & Wiper pushed through compromised patch-management servers & No \\
2014 & JPMorgan Chase & Stolen credentials on a server without two-factor authentication & No \\
2015 & Carbanak/Anunak campaign (many banks) & Spear-phishing with known Office exploits then manual operations & No \\
2015 & Tien Phong Bank (attempt) & Malware targeting SWIFT PDF reader workflow & No \\
2016 & Bangladesh Bank & Malware in SWIFT environment with stolen operator credentials & No \\
2016 & Government Savings Bank ATMs\textsuperscript{TH} & ATM malware (jackpotting) & No \\
2017 & Polish banks via KNF website & Watering hole on regulator website with known browser-plugin exploits & No (n-day) \\
2017 & National Bank of Ukraine and commercial banks (NotPetya) & Supply-chain update of accounting software plus EternalBlue worming & Supply chain \\
2018 & SPEI participants (Banco de Mexico system) & Compromise of participant-side connection software & No \\
2018 & Banco de Chile & Wiper as diversion during fraudulent SWIFT transfers & No \\
2018 & Kasikornbank and Krungthai Bank\textsuperscript{TH} & Web-application vulnerabilities in customer-facing systems & No \\
2018 & Cosmos Bank & ATM switch compromise and SWIFT fraud & No \\
2019 & Capital One & Server-side request forgery against misconfigured cloud firewall & No \\
2019 & Travelex & Ransomware via unpatched Pulse Secure VPN & No (n-day) \\
2019 & European Central Bank BIRD website & Malware injected into externally hosted website & No \\
2020 & SolarWinds Orion downstream customers & Trojanised software update (supply chain) & Supply chain \\
2020 & New Zealand stock exchange & Volumetric DDoS & No \\
2021 & Reserve Bank of New Zealand & Zero-day exploitation of Accellion FTA file-transfer appliance & Yes \\
2021 & European Banking Authority & Microsoft Exchange zero-days (ProxyLogon) & Yes \\
2021 & Thai card-not-present fraud wave\textsuperscript{TH} & Card-number enumeration and merchant-side abuse & No \\
2022 & Bank Indonesia & Ransomware (vector not disclosed) & Unknown \\
2022 & Bank of Zambia & Ransomware (vector not disclosed) & Unknown \\
2022 & Thai mobile-banking remote-access malware scams\textsuperscript{TH} & Social engineering to install remote-access malware on Android phones & No \\
2023 & Hatch Bank and others (GoAnywhere campaign) & Zero-day in Fortra GoAnywhere MFT & Yes \\
2023 & Financial firms and their service providers (MOVEit campaign) & Zero-day in MOVEit Transfer & Yes \\
2023 & ICBC Financial Services & Ransomware reportedly via Citrix Bleed after patch release & No (n-day; suspected) \\
2023 & Bank of America customers via Infosys McCamish & Ransomware at outsourcer (vector not disclosed) & Unknown \\
2023 & Latitude Financial & Stolen employee credentials at a vendor & No \\
2024 & Patelco Credit Union & Ransomware (vector not disclosed) & Unknown \\
2024 & Evolve Bank \& Trust & Ransomware after phishing & No \\
2024 & Banks and payment firms (CrowdStrike update outage) & Faulty security-agent update (non-malicious) & Common mode \\
2025 & Bybit & Compromise of multisig wallet interface supplier & Supply chain \\
\end{longtable}
}

\section{Reproducing the results}\label{app:repro}
\begin{enumerate}[leftmargin=1.6em]
\item \texttt{bash run\_all.sh} downloads the two public inputs (CVE List V5 bulk release; CISA KEV JSON) and runs \texttt{src/s01}--\texttt{s09}.
\item \texttt{s01} flattens CVE records; \texttt{s02} flags AI-credited CVEs; \texttt{s03} builds the KEV panel; \texttt{s04} produces the empirical tables and \texttt{outputs/key\_numbers.json}; \texttt{s05} summarises the curated cases; \texttt{s06} scores the Thai groups; \texttt{s07} runs the simulation; \texttt{s08} draws figures; \texttt{s09} writes the LaTeX tables.
\item \texttt{cd paper \&\& latexmk -pdf main.tex} compiles the paper.
\end{enumerate}
Results for a later snapshot will differ, because both sources are updated continuously. Seeds are fixed, so a given snapshot reproduces exactly.
